\begin{theorem}[A shifted second sphere]\hypertarget{tag-02MN}{}\label{maps:src:E.2.3}
\StatementRecord{maps:src:E.2.3}{proved}{02MN}{\SeriesRef{maps:src:E.2.2}, \eqref{maps:eq:E.9}}


Let $X'$ be a second canonical member and choose $\varphi\in
\operatorname{PGL}_2(\mathbb C)$ so that its three exceptional
values are disjoint from those of $X$.  The smooth fivefold
$P=X\times_B\varphi^*X'$ has Euler characteristic zero and
an infinite fundamental group with presentation
\begin{equation}
\begin{split}
G=\langle a_1,a_2,b_1,b_2,h_1,h_2\mid\ &
 h_1,h_2\ \text{central},\quad
 a_1^3=h_1,\ a_2^4=h_1^{-1},\\
 &b_1^3=h_2,\ b_2^4=h_2^{-1},\quad
 a_1a_2b_1b_2=1\rangle .
\end{split}
\hypertarget{tag-029M}{}\label{maps:eq:E.16}
\end{equation}
Its abelianization is $\mathbb Z$, and its rational Betti numbers are
\begin{equation}
 (1,1,1,20,63,88,63,20,1,1,1).
 \hypertarget{tag-029N}{}\label{maps:eq:E.17}
\end{equation}


\end{theorem}


\begin{proof}
\begin{enumerate}[label=\arabic*.,leftmargin=*]
\item Let $f:X\to B$ and $g:\varphi^*X'\to B$ be the two maps.
At a point $(x,x')$ of the fibre product, their disjoint critical-value
sets imply that at least one of $df_x$ and $dg_{x'}$ is surjective.
Consequently the difference differential
$(v,w)\mapsto df_x(v)-dg_{x'}(w)$ maps
$T_xX\oplus T_{x'}\varphi^*X'$ onto $T_{f(x)}B$.
In a local coordinate on $B$ it is the differential of the defining
function $f-g$. The holomorphic implicit function theorem therefore
makes $P$ smooth of complex dimension $3+3-1=5$.
It is compact because it is closed in the compact product
$X\times\varphi^*X'$.  At each special value the other
factor is a smooth four-torus.  The Euler characteristic of every
fibre is therefore zero, and Euler additivity gives $\chi(P)=0$.
We apply van Kampen with the cusp relations.
Over the six-punctured sphere the fibre group is
$\Lambda_X\oplus\Lambda_{X'}$ and the base is free on five
peripheral loops.  Each monodromy acts on its own direct summand.
Filling the first cusp kills its angular generators
$\widehat w,\widehat\delta$ and its cusp meridian.  Their normal
closure under $A_1$ also kills $\widehat u$, since
$A_1\widehat w=\widehat u-\widehat w$.  The same calculation
in the other summand leaves only $h_1=\widehat\gamma$ and
$h_2=\widehat\gamma'$.  The finite monodromies now act trivially
on them.  The four finite filling relations are precisely those
in \eqref{maps:eq:E.16}, since $\gamma(v_1)=1$ and $\gamma(v_2)=-1$.
Both cusp collars are untwisted; their meridians are killed.
The product of the remaining four meridians gives the last
relation.  These are all relations introduced by the six
fillings, so van Kampen gives exactly \eqref{maps:eq:E.16}.

\item Quotienting by $h_1,h_2$ gives the orbifold group of the sphere
with cone orders $3,4,3,4$.  Its orbifold Euler characteristic
is $-5/6$, so its universal orbifold cover is the disc and the
group is infinite.  Equivalently one can build the Seifert
$T^2$-bundle with these four primitive translation vectors; over
the disc it is a trivial torus bundle.  Its exact sequence shows
that the central subgroup in \eqref{maps:eq:E.16} is $\mathbb Z^2$.
The abelianized five-by-six relation matrix has Smith factors
$(1,1,1,1,1)$ by integral row and column operations, proving $G_{\rm ab}=\mathbb Z$.

\item For the rational homology there is no extra cusp kernel: each
cusp is the rank-two cusp of the sphere fibration times a smooth torus, and its
specialization is an isomorphism onto invariants over $\mathbb Q$.
The two monodromy groups act on independent summands.  In exterior
degrees $q=0,\ldots,8$, the $H^0$ and $H^2$ rows are both
\[
 (1,2,3,4,5,4,3,2,1),
\]
and the $H^1$ row is
\[
 (0,0,18,60,86,60,18,0,0).
\]
The local invariant polynomials of the first factor are
$A(z)=1+2z+2z^2+2z^3+z^4$ at a finite point and
$C(z)=A(z)+2z^2$ at its cusp.  Hence the Euler polynomial of
the direct-image rows is
\[
 -4(1+z)^8+(4A(z)+2C(z))(1+z)^4.
\]
Formula \eqref{maps:eq:E.9}, with six point strata, gives the displayed rows.

\item Each invariant algebra of a factor has one generator in degrees
zero through four.  Its invariant-to-coinvariant map is nonzero
in even degrees and zero in odd degrees.  In the marking used
above, this follows from the invariant form
$6\gamma\wedge\delta+u\wedge w$ and its exterior products:
the odd invariant classes pair to zero with every invariant
complementary class, and the even ones have nonzero square or
top pairing.  Its four transgressions are all nonzero, as the
factor is a sphere.  The derivation rule therefore sends a tensor
of bidegree $(i,j)$ into bidegrees $(i-1,j)$ and $(i,j-1)$,
with the first term nonzero exactly when $i>0$ and $j$ is even,
and the second exactly when $j>0$ and $i$ is even.  The nonzero
column supports are disjoint except for pairs in a single row,
so their scalar normalizations do not affect ranks.  This gives
\[
 (r_1,\ldots,r_8)=(1,2,2,3,3,2,2,1).
\]
Row reduction of these exterior-product matrices gives the displayed ranks.
There are no later differentials.  Their diagonal sums give
\eqref{maps:eq:E.17}, which also satisfies Poincar\'e duality and Euler zero.

\end{enumerate}
\end{proof}